\section{Related Work}
\label{sec:related}

\subsection{Respiratory Sound Analysis}
Public datasets have driven progress in this area. ICBHI 2017~\cite{rocha2019open} provides 920 recordings from 126 subjects with respiratory-cycle and crackle/wheeze labels. SPRSound~\cite{zhang2022sprsound} is the first open paediatric database with event-level onset/offset labels, and it is used in the IEEE BioCAS grand challenges~\cite{zhang2022grand,zhang2024meta}. HF\_Lung~\cite{hsu2021benchmarking,hsu2022progressively} adds breath-phase and adventitious-sound annotations, and DeepBreath~\cite{heitmann2023deepbreath} aggregates multi-site recordings into patient-level predictions.

On these datasets, convolutional networks on spectrograms~\cite{bardou2018lung,demir2020convolutional,demir2020classification} and hybrid convolutional neural network (CNN)--recurrent models~\cite{acharya2020deep,petmezas2022automated} are now standard. Accuracy has been pushed further by attention-based feature fusion~\cite{crisdayanti2024attention,dong2025respiratory,kim2025enhanced}, large pretrained audio models~\cite{kong2020panns,gong2021ast}, and respiratory-specific pretraining, augmentation and distillation~\cite{bae2023patchmix,niizumi2025towards,toikkanen2025improving,castejonbarrio2026leveraging,hirasawa2026domain}. Other work improves robustness to patients, devices and noise~\cite{kim2025adaptive,jeong2025patient,tzeng2025improving}. The strength of this line of work is accuracy on curated benchmarks. Its weakness, for our purpose, is that almost all of it is offline: models see whole recordings or pre-segmented cycles and spend uniform computation on every frame. Acharya and Basu~\cite{acharya2020deep} target wearables, but through weight quantization and patient-specific tuning, not by adapting computation to the signal.

\subsection{Respiratory Event Detection}
Event detection locates the onset and offset of each respiratory event in addition to its class. Frame-wise recurrent detectors~\cite{hsu2021benchmarking} and convolutional-recurrent networks~\cite{cakir2017convolutional} are the common backbones. EZhouNet~\cite{chu2026ezhounet} uses a graph neural network with anchor intervals on SPRSound and reaches strong event-level scores. Kala and Elhilali~\cite{kala2026multistage} first select anomaly-specific segments and then classify them. These methods show that graph models and segment selection help detection accuracy. However, they select \emph{which segments to classify} offline, not \emph{which parts of a running state to recompute} under a compute budget. Detection is usually scored with event- and segment-based metrics~\cite{mesaros2016metrics} using tolerance collars in the style of the Detection and Classification of Acoustic Scenes and Events (DCASE) challenges~\cite{turpault2019sound}; threshold-free alternatives also exist~\cite{bilen2020framework,ebbers2022threshold}.

\subsection{Conditional Computation}
\label{sec:related-cc}
Several methods reduce computation by skipping work that does not seem needed. Delta networks~\cite{neil2017delta} transmit recurrent activations only when they change beyond a threshold. This is simple and hardware-friendly, but it equates importance with change magnitude. Skip RNN~\cite{campos2018skip} learns a gate that skips whole state updates under a budget penalty. Variable-computation recurrent neural networks (RNNs)~\cite{jernite2017variable} and Skim-RNN~\cite{seo2018neural} update only part of the hidden state per step. Adaptive computation time~\cite{graves2016adaptive} learns how many internal steps to take, and DynamicViT~\cite{rao2021dynamicvit} learns which transformer tokens to drop; see~\cite{han2022dynamic} for a survey. These learned gates are trained end to end through relaxations or penalties, so the budget is met only on average and the gate never sees an explicit estimate of what skipping costs. Our controller instead regresses a counterfactual estimate of each skip's effect on the output, and a selection rule enforces the budget exactly at every step.

\subsection{Temporal Graph Networks and AGN-Net}
Temporal graph networks combine message passing, often with graph attention~\cite{velickovic2018graph}, with node states that evolve over time. GCRN~\cite{seo2018structured} places graph convolutions inside a recurrent cell. EvolveGCN~\cite{pareja2020evolvegcn} evolves the weights of a graph convolutional network (GCN) with an RNN. TGN~\cite{rossi2020temporal} updates node memories only for nodes involved in an event. These models capture structure and memory well, but their update sets are fixed by the model or by the data, not chosen to save computation. AGN-Net~\cite{islam2026agnnet} is the exception. It scores nodes with a learned priority, updates only a $k$-hop region around high-priority nodes, and carries the remaining states forward with a temporal-momentum update. Its priority is trained to predict next-step representation change, and the study reported that this priority did not track realized change and improved on random selection only modestly; most of its benefit came from the auxiliary loss acting as a regularizer. We keep AGN-Net's localized update mechanism but replace its training target with the task-level consequence of skipping, which we show the controller can learn and exploit.
